% !TeX root = ../manual.tex
%=======================================================================
%  config/plantilla.tex — mecanismos propios de este manual
%=======================================================================

%----------------------- Inclusión de capturas --------------------------
% \captura{archivo}{ancho}{pie}{etiqueta} — figura centrada con la captura
% de pantalla, tomada de figuras/capturas/. Si el archivo no existe, deja
% un recuadro visible en vez de fallar la compilación.
\newcommand{\faltacaptura}[1]{%
  \fbox{\parbox{0.92\linewidth}{\centering\small
    \textcolor{BrickRed}{\textbf{Falta \texttt{figuras/capturas/#1.jpg}}}}}}

\newcommand{\captura}[4]{%
  \begin{figure}[H]
    \centering
    \IfFileExists{figuras/capturas/#1.jpg}%
      {\includegraphics[width=#2\linewidth]{capturas/#1}}%
      {\faltacaptura{#1}}
    \caption{#3}
    \label{fig:#4}
  \end{figure}}

%--------------------------- Cajas de aviso -----------------------------
\definecolor{notaback}{HTML}{EFF6FF}
\definecolor{notaframe}{HTML}{1168BD}
\definecolor{avisoback}{HTML}{FEF3C7}
\definecolor{avisoframe}{HTML}{B45309}

% \nota{...} — información útil o aclaración, sin urgencia.
\newcommand{\nota}[1]{%
  \begin{tcolorbox}[enhanced,breakable,
      colback=notaback,colframe=notaframe,boxrule=0.6pt,arc=2pt,
      left=8pt,right=8pt,top=5pt,bottom=5pt,
      title={Nota},fonttitle=\bfseries\small,
      coltitle=notaframe,colbacktitle=notaback,
      fontupper=\small]
    #1
  \end{tcolorbox}}

% \aviso{...} — advertencia sobre algo irreversible o que requiere cuidado.
\newcommand{\aviso}[1]{%
  \begin{tcolorbox}[enhanced,breakable,
      colback=avisoback,colframe=avisoframe,boxrule=0.6pt,arc=2pt,
      left=8pt,right=8pt,top=5pt,bottom=5pt,
      title={Atención},fonttitle=\bfseries\small,
      coltitle=avisoframe,colbacktitle=avisoback,
      fontupper=\small]
    #1
  \end{tcolorbox}}

%------------------------- Referencias a la interfaz ---------------------
% \boton{Guardar} — nombre de un botón, campo o elemento de la interfaz.
\newcommand{\boton}[1]{\textbf{#1}}
% \ruta{Canchas > Editar} — ruta de navegación dentro del menú.
\newcommand{\ruta}[1]{\textsf{#1}}
% \campo{usuario} — nombre de un campo de formulario.
\newcommand{\campo}[1]{\textit{#1}}
